\documentclass[UTF8,12pt,a4paper,fontset=fandol]{ctexart}

% 支持从项目根目录或当前笔记目录编译。
\makeatletter
\def\input@path{{../}}
\makeatother
\usepackage{researchnotes}

\title{双阶乘笔记}
\author{}
\date{}

\begin{document}
\maketitle

\section{定义与基本计算}

普通阶乘 $n!$ 将从 $1$ 到正整数 $n$ 的所有整数相乘。与之类似，
\keyterm{双阶乘}（double factorial）记作 $n!!$，但相邻因子每次相差 $2$。

\begin{definition}[双阶乘]
对于正整数 $n$，定义
\[
  n!!=n(n-2)(n-4)\cdots,
\]
其中，$n$ 为奇数时最后一个因子为 $1$，$n$ 为偶数时最后一个因子为 $2$。
此外，约定 $0!!=1$ 和 $(-1)!!=1$。
\end{definition}

将奇数与偶数分别表示为 $2m-1$ 和 $2m$，对于整数 $m\geq 1$，有
\begin{align*}
  (2m-1)!!&=1\times3\times5\times\cdots\times(2m-1),\\
  (2m)!!&=2\times4\times6\times\cdots\times2m.
\end{align*}
上述约定使递推关系 $n!!=n(n-2)!!$ 对所有整数 $n\geq1$ 均成立，
包括 $1!!=1\cdot(-1)!!=1$ 和 $2!!=2\cdot0!!=2$。
这里的 $(-1)!!$ 是约定的值，并不表示对所有负整数都已定义双阶乘。

\begin{example}[奇数与偶数的双阶乘]
计算时，每次将因子减去 $2$，直到到达相应的终点：
\[
  7!!=7\times5\times3\times1=105,\qquad
  8!!=8\times6\times4\times2=384.
\]
\end{example}

\section{与普通阶乘的关系}

对于整数 $m\geq1$，从偶数双阶乘的每个因子中提出一个 $2$，可得
\[
  (2m)!!=(2\cdot1)(2\cdot2)\cdots(2\cdot m)
  =2^m(1\cdot2\cdots m).
\]
因此，
\begin{equation}\label{eq:even-double-factorial}
  (2m)!!=2^m m!.
\end{equation}
例如，$8!!=2^4\cdot4!=16\times24=384$，与直接计算的结果一致。

将 $(2m)!$ 中的因子按奇偶性分组，恰好得到两个双阶乘：
\begin{equation}\label{eq:factorial-split}
  (2m)!=(2m)!!(2m-1)!!.
\end{equation}
这是因为 $1$ 到 $2m$ 中的所有偶数构成 $(2m)!!$，所有奇数构成
$(2m-1)!!$，每个因子都恰好出现一次。
由式~\eqref{eq:even-double-factorial} 和式~\eqref{eq:factorial-split} 可得
\begin{equation}\label{eq:odd-double-factorial}
  (2m-1)!!=\frac{(2m)!}{(2m)!!}=\frac{(2m)!}{2^m m!}.
\end{equation}
例如，$7!!=8!/(2^4\cdot4!)=105$。
结合 $0!=0!!=(-1)!!=1$ 的约定，上述三个等式在 $m=0$ 时也成立。

\section{记号辨析与应用背景}

记号 $n!!$ 表示双阶乘，\keyterm{不能理解为对阶乘的结果再取一次阶乘}。
例如，
\[
  5!!=5\times3\times1=15,\qquad (5!)!=120!.
\]
前者按步长 $2$ 连乘，后者先计算 $5!=120$，再计算 $120$ 的普通阶乘，
两者表示不同的运算。

双阶乘常见于组合数学、概率论、高斯积分、球的体积与球面的面积计算，
以及某些函数的泰勒展开。这些问题中常出现间隔为 $2$ 的连乘积，
双阶乘为它们提供了简洁的记号；利用它与普通阶乘的关系，还可以进一步化简计算。

\end{document}
